\documentclass[a4paper, 14pt]{extarticle}

% Настройка поддержки русского языка и системных шрифтов
\usepackage{fontspec}
\usepackage[russian]{babel}

% Установка заданных шрифтов
\setmainfont{Times New Roman}
\setmonofont{Courier New}

% Настройка полей страницы
\usepackage{geometry}
\geometry{a4paper, top=2cm, bottom=2cm, left=3cm, right=1.5cm}

% Пакет для оформления сниппетов кода
\usepackage{listings}
\usepackage{xcolor}

\lstset{
    basicstyle=\ttfamily\small, % Использование моноширинного шрифта (Courier New)
    frame=single,
    breaklines=true,
    backgroundcolor=\color{gray!5},
    xleftmargin=0.5cm,
    showstringspaces=false
}

% Пакет для кликабельных ссылок
\usepackage[hidelinks]{hyperref}

\begin{document}

\title{Описание репозитория \texttt{texlive-container}}
\date{}
\maketitle

\section*{Обзор проекта}

Репозиторий 
\href{https://github.com/denisstrizhkin/texlive-container}{texlive-container} 
представляет собой проект конфигурации для развертывания
издательской системы \TeX\ Live в изолированном контейнере. 

Ключевой особенностью данного решения является использование дистрибутива
Alpine Linux в качестве базового образа системы.
Это позволяет получить легковесную и готовую к работе среду для компиляции 
\LaTeX-документов без необходимости устанавливать
множество дополнительных пакетов и зависимостей в основную операционную систему.

\section*{Использование и сборка образа}

Для работы с проектом и локальной сборки образа контейнера используется
инструмент \texttt{podman} (или \texttt{docker}). В официальной документации репозитория 
приводится следующая команда для загрузки исходного кода
и запуска процесса сборки:

\begin{lstlisting}[language=bash]
podman build -t texlive https://github.com/denisstrizhkin/texlive-container.git
\end{lstlisting}

После успешного выполнения данной команды в вашей системе появится готовый
образ с именем \texttt{texlive}.
Его можно использовать для автоматизации процессов сборки научных статей,
отчетов и презентаций, а также для интеграции в CI/CD пайплайны.

Чтобы скомпилировать документ внутри контейнера, необходимо примонтировать 
текущую рабочую директорию. Например, данный документ можно скомпилировать следующей командой:
\begin{lstlisting}[language=bash]
podman run --rm -v "$PWD":/data -w /data texlive lualatex sample.tex
\end{lstlisting}

При использовании \texttt{docker} команда аналогична:
\begin{lstlisting}[language=bash]
docker run --rm -v "$PWD":/data -w /data texlive lualatex sample.tex
\end{lstlisting}

\end{document}
